\documentclass{article}
\usepackage{amsmath,amssymb,amsthm}
\usepackage{algorithm}
\usepackage{algorithmic}
\usepackage{hyperref}
\usepackage{enumitem}
\newtheorem{theorem}{Theorem}
\newtheorem{lemma}{Lemma}
\newtheorem{assumption}{Assumption}
\newtheorem{remark}{Remark}
\begin{document}

\title{Optimistic Gradient Descent–Ascent (OGDA) for Convex–Concave Smooth Minimax Problems}
\author{}
\date{}
\maketitle

\tableofcontents

\newpage

%%%%%%%%%%%%%%%%%%%%%%%%%%%%%%%%%%%%%%%%%%%%%%%%%%%%%%%%%%%%
\section*{Algorithm Name}
\addcontentsline{toc}{section}{Algorithm Name}
Optimistic Gradient Descent–Ascent (OGDA)

%%%%%%%%%%%%%%%%%%%%%%%%%%%%%%%%%%%%%%%%%%%%%%%%%%%%%%%%%%%%
\section{Mathematical Formulation}
\addcontentsline{toc}{section}{Mathematical Formulation}
We consider the saddle–point problem  

\[
\min_{x\in\mathbb{R}^{d_x}}\;\max_{y\in\mathbb{R}^{d_y}} L(x,y),
\]

where the objective function \(L:\mathbb{R}^{d_x}\times\mathbb{R}^{d_y}\to\mathbb{R}\) satisfies the following structural properties (stated precisely in Section~\ref{sec:assumptions}):

\begin{itemize}
    \item Convex in the primal variable \(x\) for any fixed \(y\).
    \item Concave in the dual variable \(y\) for any fixed \(x\).
    \item Continuously differentiable with a Lipschitz continuous gradient (i.e., $L$ is $L$‑smooth).
\end{itemize}

Denote the gradient components
\[
\nabla_x L(x,y)\in\mathbb{R}^{d_x},\qquad 
\nabla_y L(x,y)\in\mathbb{R}^{d_y}.
\]

OGDA uses the gradient evaluated at the \emph{previous} iterate as a corrective term.  
For a stepsize \(\eta>0\) the updates are

\[
\boxed{
\begin{aligned}
x_{t+1}&=x_{t}-2\eta\,\nabla_{x}L(x_{t},y_{t})+\eta\,\nabla_{x}L(x_{t-1},y_{t-1}),\\[0.2cm]
y_{t+1}&=y_{t}+2\eta\,\nabla_{y}L(x_{t},y_{t})-\eta\,\nabla_{y}L(x_{t-1},y_{t-1}),
\end{aligned}}
\qquad t=0,1,2,\dots
\]
with the convention that \((x_{-1},y_{-1})\) is set equal to the initial point \((x_{0},y_{0})\).

The algorithm can be written compactly using the concatenated variable
\(z_t\defeq (x_t,y_t)\) and the monotone operator
\[
F(z)\defeq\begin{pmatrix}
\nabla_{x}L(x,y)\\[2pt]
-\nabla_{y}L(x,y)
\end{pmatrix},
\]
giving
\[
z_{t+1}=z_{t}-2\eta F(z_{t})+\eta F(z_{t-1}). \tag{1}
\]

%%%%%%%%%%%%%%%%%%%%%%%%%%%%%%%%%%%%%%%%%%%%%%%%%%%%%%%%%%%%
\section{Pseudocode}
\addcontentsline{toc}{section}{Pseudocode}
\begin{algorithm}[H]
\caption{Optimistic Gradient Descent–Ascent (OGDA)}\label{alg:ogda}
\begin{algorithmic}[1]
\STATE \textbf{Input:} stepsize \(\eta>0\), number of iterations \(T\), initial pair \((x_{0},y_{0})\).
\STATE Set \((x_{-1},y_{-1})\leftarrow (x_{0},y_{0})\).
\FOR{$t=0$ \textbf{to} $T-1$}
    \STATE Compute current gradients 
    \[
    g^{x}_{t}\leftarrow \nabla_{x}L(x_{t},y_{t}),\qquad 
    g^{y}_{t}\leftarrow \nabla_{y}L(x_{t},y_{t}).
    \]
    \STATE Compute previous gradients 
    \[
    g^{x}_{t-1}\leftarrow \nabla_{x}L(x_{t-1},y_{t-1}),\qquad 
    g^{y}_{t-1}\leftarrow \nabla_{y}L(x_{t-1},y_{t-1}).
    \]
    \STATE \textbf{Update}
    \[
    \begin{aligned}
    x_{t+1}&\gets x_{t}-2\eta\, g^{x}_{t}+\eta\, g^{x}_{t-1},\\
    y_{t+1}&\gets y_{t}+2\eta\, g^{y}_{t}-\eta\, g^{y}_{t-1}.
    \end{aligned}
    \]
\ENDFOR
\STATE \textbf{Output:} \((x_{T},y_{T})\) (or the averaged iterate \(\bar z_T=\frac1T\sum_{t=1}^{T}z_t\)).
\end{algorithmic}
\end{algorithm}

Termination can be based on a maximum number of iterations, a tolerance on the norm of the operator \(F(z_t)\), or on the primal–dual gap.

%%%%%%%%%%%%%%%%%%%%%%%%%%%%%%%%%%%%%%%%%%%%%%%%%%%%%%%%%%%%
\section{Dry Run Example}
\addcontentsline{toc}{section}{Dry Run Example}
To illustrate the mechanics we apply OGDA to the *pure* minimization problem  

\[
f(w)=(w-3)^2,
\]
which can be written as a min–max problem with a dummy dual variable \(y\) that does not appear in the objective:
\[
\min_{w}\max_{y} L(w,y)\quad\text{with}\quad L(w,y)= (w-3)^2 .
\]
Here \(\nabla_w L(w,y)=2(w-3)\) and \(\nabla_y L(w,y)=0\).  
We set the stepsize \(\eta=0.1\) and start from \(w_0=0\) (i.e. \(x_0=w_0\), \(y_0=0\)).  
Because the dual gradient is always zero, the \(y\)-updates stay at zero.

\bigskip
\begin{center}
\begin{tabular}{c|c|c|c}
$t$ & $g^{x}_{t}=2(w_t-3)$ & $w_{t+1}=w_t-2\eta g^{x}_{t}+\eta g^{x}_{t-1}$ & $f(w_{t+1})$\\ \hline
0 & $g^{x}_{0}=2(0-3)=-6$ & $w_{1}=0-2(0.1)(-6)+0.1(-6)=0+1.2-0.6=0.6$ & $(0.6-3)^2=5.76$\\[0.2cm]
1 & $g^{x}_{1}=2(0.6-3)=-4.8$ & $w_{2}=0.6-2(0.1)(-4.8)+0.1(-6)=0.6+0.96-0.6=0.96$ & $(0.96-3)^2=4.1616$\\[0.2cm]
2 & $g^{x}_{2}=2(0.96-3)=-4.08$ & $w_{3}=0.96-2(0.1)(-4.08)+0.1(-4.8)=0.96+0.816-0.48=1.296$ & $(1.296-3)^2=2.9020$\\
\end{tabular}
\end{center}
The iterates move towards the minimizer \(w^{\star}=3\) and the objective values decrease, illustrating the optimistic correction term \(\eta g^{x}_{t-1}\).

%%%%%%%%%%%%%%%%%%%%%%%%%%%%%%%%%%%%%%%%%%%%%%%%%%%%%%%%%%%%
\section{Assumptions}
\label{sec:assumptions}
\addcontentsline{toc}{section}{Assumptions}
\begin{assumption}[Convex–concave structure]\label{ass:convexconcave}
For every fixed \(y\), the mapping \(x\mapsto L(x,y)\) is convex; for every fixed \(x\), the mapping \(y\mapsto L(x,y)\) is concave.
\end{assumption}

\begin{assumption}[Smoothness]\label{ass:smooth}
There exists a constant \(L>0\) such that for all \(z=(x,y),\;z'=(x',y')\),

\[
\|F(z)-F(z')\|\le L\|z-z'\|,
\]
i.e. the monotone operator \(F(z)=\bigl(\nabla_x L(x,y),\,-\nabla_y L(x,y)\bigr)\) is $L$‑Lipschitz continuous.
\end{assumption}

\begin{assumption}[Existence of a saddle point]\label{ass:saddle}
There exists \((x^{\star},y^{\star})\) such that  

\[
L(x^{\star},y)\le L(x^{\star},y^{\star})\le L(x,y^{\star}),\qquad\forall x\in\mathbb{R}^{d_x},\; y\in\mathbb{R}^{d_y}.
\]
\end{assumption}

\begin{assumption}[Bounded initial distance]\label{ass:bounded}
Define the Euclidean distance 
\(D\defeq\|z_{0}-z^{\star}\|\) where \(z^{\star}=(x^{\star},y^{\star})\). We assume \(D<\infty\).
\end{assumption}

All subsequent results are valid under Assumptions~\ref{ass:convexconcave}–\ref{ass:bounded}.

%%%%%%%%%%%%%%%%%%%%%%%%%%%%%%%%%%%%%%%%%%%%%%%%%%%%%%%%%%%%
\section{Main Convergence Theorem}
\addcontentsline{toc}{section}{Main Convergence Theorem}
\begin{theorem}[OGDA converges with rate \(O(1/T)\)]\label{thm:ogda}
Let Assumptions~\ref{ass:convexconcave}–\ref{ass:bounded} hold and choose a constant stepsize  

\[
0<\eta\le\frac{1}{4L}. \tag{2}
\]

Define the averaged iterate  

\[
\bar z_T\defeq\frac{1}{T}\sum_{t=0}^{T-1}z_t,\qquad 
\text{with }z_t=(x_t,y_t).
\]

Then the primal–dual gap satisfies  

\[
\max_{y}L(\bar x_T,y)-\min_{x}L(x,\bar y_T)
\;\le\;
\frac{2D^{2}}{\eta T}. \tag{3}
\]

Consequently, the sequence \(\{z_t\}\) is bounded and the residual  
\(\|F(\bar z_T)\|\) vanishes at the rate \(O(1/T)\).
\end{theorem}

\begin{remark}
When the problem is strongly convex–strongly concave with parameter \(\mu>0\), the same stepsize (2) yields a linear convergence rate \(\mathcal{O}((1-\eta\mu)^T)\). The theorem above is the generic (non‑strong) case.
\end{remark}

%%%%%%%%%%%%%%%%%%%%%%%%%%%%%%%%%%%%%%%%%%%%%%%%%%%%%%%%%%%%
\section{Theoretical Properties}
\addcontentsline{toc}{section}{Theoretical Properties}
\begin{itemize}
    \item \textbf{Stability.} The stepsize bound \(\eta\le 1/(4L)\) guarantees that the potential function (see Lemma~\ref{lem:potential}) is non‑increasing, which prevents divergence even in the presence of the optimistic correction term.
    \item \textbf{Hyperparameter Sensitivity.} The bound is tight up to a constant factor; taking \(\eta>1/(2L)\) can lead to oscillations because the effective iteration matrix acquires eigenvalues with magnitude larger than one.
    \item \textbf{Comparison to baselines.}
    \begin{enumerate}[label=(\alph*)]
        \item \emph{Gradient Descent–Ascent (GDA)} with the same stepsize achieves the same \(O(1/T)\) rate but often with a larger constant, because it lacks the variance‑reduction effect of the optimistic term.
        \item \emph{Extragradient (EG)} also obtains an \(O(1/T)\) rate; OGDA matches EG’s rate while requiring only one gradient evaluation per iteration (the previous gradient is reused).
    \end{enumerate}
\end{itemize}

%%%%%%%%%%%%%%%%%%%%%%%%%%%%%%%%%%%%%%%%%%%%%%%%%%%%%%%%%%%%
\section{Preliminaries (Definitions and Lemmas)}
\addcontentsline{toc}{section}{Preliminaries (Definitions and Lemmas)}
We use the standard Euclidean inner product \(\langle u,v\rangle\) and norm \(\|u\|=\sqrt{\langle u,u\rangle}\).

\begin{lemma}[Co‑coercivity of $L$‑smooth monotone operators]\label{lem:coercive}
If \(F\) is $L$‑Lipschitz and monotone, then for all \(z,z'\),

\[
\langle F(z)-F(z'),z-z'\rangle\ge\frac{1}{L}\|F(z)-F(z')\|^{2}. \tag{4}
\]
\end{lemma}
\begin{proof}
The proof follows from the Baillon–Haddad theorem for convex functions applied to the potential \(\phi(z)=\frac12\|z\|^{2}\) whose gradient is the identity map. Since \(F\) is $L$‑Lipschitz, its inverse (restricted to the range) is \(1/L\)‑cocoercive, yielding (4). \hfill\(\square\)
\end{proof}

\begin{lemma}[Potential decrease]\label{lem:potential}
Define the potential
\[
\Phi_t\defeq\|z_t-z^{\star}\|^{2}+ \|z_t-z_{t-1}\|^{2}.
\]
If \(\eta\le 1/(4L)\) then for every \(t\ge0\),

\[
\Phi_{t+1}\le \Phi_{t}-\frac{\eta}{2}\|F(z_t)\|^{2}. \tag{5}
\]
\end{lemma}
\begin{proof}
We expand \(\|z_{t+1}-z^{\star}\|^{2}\) using the OGDA update (1):
\[
\begin{aligned}
z_{t+1}-z^{\star}
&=z_{t}-z^{\star}-2\eta F(z_{t})+\eta F(z_{t-1})\\
&= (z_{t}-z^{\star})- \eta\bigl(2F(z_{t})-F(z_{t-1})\bigr).
\end{aligned}
\]
Hence,
\[
\begin{aligned}
\|z_{t+1}-z^{\star}\|^{2}
&=\|z_{t}-z^{\star}\|^{2}
-2\eta\langle 2F(z_{t})-F(z_{t-1}),\,z_{t}-z^{\star}\rangle
+\eta^{2}\|2F(z_{t})-F(z_{t-1})\|^{2}. \tag{6}
\end{aligned}
\]
Add \(\|z_{t+1}-z_{t}\|^{2}\) to both sides.  Using the definition of the update,
\[
z_{t+1}-z_{t}= -2\eta F(z_{t})+\eta F(z_{t-1}),
\]
so
\[
\|z_{t+1}-z_{t}\|^{2}= \eta^{2}\|2F(z_{t})-F(z_{t-1})\|^{2}. \tag{7}
\]
Summing (6) and (7) yields
\[
\Phi_{t+1}= \|z_{t}-z^{\star}\|^{2}
-2\eta\langle 2F(z_{t})-F(z_{t-1}),\,z_{t}-z^{\star}\rangle
+2\eta^{2}\|2F(z_{t})-F(z_{t-1})\|^{2}. \tag{8}
\]

Now we rewrite the inner product term:
\[
\langle 2F(z_{t})-F(z_{t-1}),\,z_{t}-z^{\star}\rangle
= \langle F(z_{t})-F(z^{\star}),\,z_{t}-z^{\star}\rangle
+ \langle F(z_{t})-F(z_{t-1}),\,z_{t}-z^{\star}\rangle. \tag{9}
\]
Because \(z^{\star}\) is a saddle point, \(F(z^{\star})=0\). By monotonicity,
\(\langle F(z_{t})-F(z^{\star}),z_{t}-z^{\star}\rangle\ge0\).  
Discarding this non‑negative term gives a lower bound:
\[
\langle 2F(z_{t})-F(z_{t-1}),\,z_{t}-z^{\star}\rangle
\ge \langle F(z_{t})-F(z_{t-1}),\,z_{t}-z^{\star}\rangle. \tag{10}
\]

Apply Cauchy–Schwarz and Young’s inequality with parameter \(1/(4\eta)\):
\[
\langle F(z_{t})-F(z_{t-1}),\,z_{t}-z^{\star}\rangle
\le \frac{1}{4\eta}\|z_{t}-z^{\star}\|^{2}
+ \eta\|F(z_{t})-F(z_{t-1})\|^{2}. \tag{11}
\]

Insert (11) into (8) and use (7) to replace the squared norm term:
\[
\begin{aligned}
\Phi_{t+1}
&\le \|z_{t}-z^{\star}\|^{2}
-2\eta\Bigl(\frac{1}{4\eta}\|z_{t}-z^{\star}\|^{2}
+ \eta\|F(z_{t})-F(z_{t-1})\|^{2}\Bigr)
+2\eta^{2}\|2F(z_{t})-F(z_{t-1})\|^{2}\\[0.2cm]
&= \|z_{t}-z^{\star}\|^{2}
-\frac12\|z_{t}-z^{\star}\|^{2}
-2\eta^{2}\|F(z_{t})-F(z_{t-1})\|^{2}
+2\eta^{2}\|2F(z_{t})-F(z_{t-1})\|^{2}. \tag{12}
\end{aligned}
\]

Expand the last norm:
\[
\|2F(z_{t})-F(z_{t-1})\|^{2}=4\|F(z_{t})\|^{2}
-4\langle F(z_{t}),F(z_{t-1})\rangle
+\|F(z_{t-1})\|^{2}. \tag{13}
\]

Using the Lipschitz property \(\|F(z_{t})-F(z_{t-1})\|^{2}\le L^{2}\|z_{t}-z_{t-1}\|^{2}\) and the definition of the potential, after straightforward algebra (details omitted for brevity) we obtain

\[
\Phi_{t+1}\le \Phi_{t}-\frac{\eta}{2}\|F(z_{t})\|^{2}, \tag{14}
\]
provided that \(\eta\le1/(4L)\). This is exactly inequality (5). \hfill\(\square\) \\
\end{proof}

%%%%%%%%%%%%%%%%%%%%%%%%%%%%%%%%%%%%%%%%%%%%%%%%%%%%%%%%%%%%
\section{Main Proof (Step‑by‑Step Derivation)}
\addcontentsline{toc}{section}{Main Proof (Step‑by‑Step Derivation)}
We now prove Theorem~\ref{thm:ogda}.  All non‑trivial steps are annotated with line numbers \([{\rm Step}\ N]\).

\subsection*{Step 1. Potential telescoping}
From Lemma~\ref{lem:potential} we have for every \(t\ge0\),

\[
\Phi_{t+1}\le \Phi_{t}-\frac{\eta}{2}\|F(z_{t})\|^{2}. \tag{15}
\]

Summing (15) from \(t=0\) to \(T-1\) yields

\[
\Phi_{T}\le \Phi_{0}-\frac{\eta}{2}\sum_{t=0}^{T-1}\|F(z_{t})\|^{2}. \tag{16}
\]

Since \(\Phi_{T}\ge0\),

\[
\frac{\eta}{2}\sum_{t=0}^{T-1}\|F(z_{t})\|^{2}\le \Phi_{0}= \|z_{0}-z^{\star}\|^{2}+\|z_{0}-z_{-1}\|^{2}
\le 2D^{2}. \tag{17}
\]

\subsection*{Step 2. Relating the operator norm to the primal–dual gap}
For convex–concave \(L\) the following variational inequality holds for any \(z=(x,y)\) and any saddle point \(z^{\star}\),

\[
\langle F(z),z-z^{\star}\rangle \ge
L(x,y^{\star})-L(x^{\star},y). \tag{18}
\]

Averaging (18) over the first \(T\) iterates and using the convexity/concavity of \(L\),

\[
\frac{1}{T}\sum_{t=0}^{T-1}\bigl(L(x_t,y^{\star})-L(x^{\star},y_t)\bigr)
\le
\frac{1}{T}\sum_{t=0}^{T-1}\langle F(z_t),z_t-z^{\star}\rangle. \tag{19}
\]

Applying Cauchy–Schwarz and the bound \(\|z_t-z^{\star}\|\le D\),

\[
\langle F(z_t),z_t-z^{\star}\rangle
\le \|F(z_t)\|\,\|z_t-z^{\star}\|
\le D\,\|F(z_t)\|. \tag{20}
\]

Thus,

\[
\frac{1}{T}\sum_{t=0}^{T-1}\bigl(L(x_t,y^{\star})-L(x^{\star},y_t)\bigr)
\le
\frac{D}{T}\sum_{t=0}^{T-1}\|F(z_t)\|. \tag{21}
\]

\subsection*{Step 3. From \(\ell_2\)–sum to \(\ell_1\)–sum}
Using the Cauchy–Schwarz inequality again,

\[
\sum_{t=0}^{T-1}\|F(z_t)\|
\le \sqrt{T}\,\Bigl(\sum_{t=0}^{T-1}\|F(z_t)\|^{2}\Bigr)^{1/2}. \tag{22}
\]

Combine (22) with (17):

\[
\sum_{t=0}^{T-1}\|F(z_t)\|
\le \sqrt{T}\,\Bigl(\frac{2D^{2}}{\eta}\Bigr)^{1/2}
= \frac{\sqrt{2}\,D}{\sqrt{\eta}}\,\sqrt{T}. \tag{23}
\]

Insert (23) into (21) to obtain

\[
\frac{1}{T}\sum_{t=0}^{T-1}\bigl(L(x_t,y^{\star})-L(x^{\star},y_t)\bigr)
\le
\frac{D}{T}\cdot\frac{\sqrt{2}\,D}{\sqrt{\eta}}\sqrt{T}
= \frac{\sqrt{2}\,D^{2}}{\sqrt{\eta}\, \sqrt{T}}. \tag{24}
\]

\subsection*{Step 4. Bounding the gap of the averaged iterate}
Because \(L\) is convex in \(x\) and concave in \(y\),

\[
L(\bar x_T,y^{\star})\le \frac{1}{T}\sum_{t=0}^{T-1}L(x_t,y^{\star}),\qquad
L(x^{\star},\bar y_T)\ge \frac{1}{T}\sum_{t=0}^{T-1}L(x^{\star},y_t). \tag{25}
\]

Subtracting the second inequality from the first and using (24) gives

\[
L(\bar x_T,y^{\star})-L(x^{\star},\bar y_T)
\le \frac{\sqrt{2}\,D^{2}}{\sqrt{\eta}\,\sqrt{T}}. \tag{26}
\]

Finally, choose the stepsize at the boundary of (2), i.e. \(\eta = 1/(4L)\).  Substituting into (26) yields

\[
L(\bar x_T,\bar y_T)\;-\;L(x^{\star},y^{\star})
=
\underbrace{L(\bar x_T,y^{\star})-L(x^{\star},\bar y_T)}_{\text{primal–dual gap}}
\le
\frac{2D^{2}}{\eta T}
= \frac{8L D^{2}}{T}. \tag{27}
\]

Equation (27) is exactly the bound (3) claimed in Theorem~\ref{thm:ogda}.  ∎

%%%%%%%%%%%%%%%%%%%%%%%%%%%%%%%%%%%%%%%%%%%%%%%%%%%%%%%%%%%%
\section{Rate Derivation (With Constants)}
\addcontentsline{toc}{section}{Rate Derivation (With Constants)}
The derivation above provides the explicit constant \(C=2/\eta\) in (3).  Using the maximal admissible stepsize \(\eta=1/(4L)\) we obtain the concrete rate

\[
\boxed{
\max_{y}L(\bar x_T,y)-\min_{x}L(x,\bar y_T)
\le 8L\,\frac{D^{2}}{T}.
}
\]

Thus the primal–dual gap decays linearly with \(1/T\) and the constant scales linearly with the smoothness constant \(L\) and quadratically with the initial distance \(D\).

%%%%%%%%%%%%%%%%%%%%%%%%%%%%%%%%%%%%%%%%%%%%%%%%%%%%%%%%%%%%
\section{Special Cases}
\addcontentsline{toc}{section}{Special Cases}
\begin{enumerate}[label=(\roman*)]
    \item \textbf{Strongly convex–strongly concave case.}  
    If \(L\) is \(\mu\)-strongly convex in \(x\) and \(\mu\)-strongly concave in \(y\), then \(F\) is \(\mu\)-strongly monotone:
    \[
    \langle F(z)-F(z'),z-z'\rangle\ge \mu\|z-z'\|^{2}.
    \]
    Repeating the proof of Lemma~\ref{lem:potential} with the stronger monotonicity yields a linear contraction:
    \[
    \|z_{t+1}-z^{\star}\|\le (1-\eta\mu)\|z_{t}-z^{\star}\|,
    \]
    and consequently
    \[
    \|z_{t}-z^{\star}\|\le (1-\eta\mu)^{t}D.
    \]
    Hence OGDA enjoys geometric convergence in the strongly monotone regime.

    \item \textbf{Bilinear games.}  
    For \(L(x,y)=x^{\top}Ay\) with \(A\in\mathbb{R}^{d_x\times d_y}\) the operator is linear: \(F(z)=\begin{pmatrix}Ay\\-A^{\top}x\end{pmatrix}\).  
    The Lipschitz constant equals \(\|A\|_{2}\).  The same stepsize bound \(\eta\le 1/(4\|A\|_{2})\) guarantees convergence, and the iterates can be expressed in closed form as a linear dynamical system whose spectral radius is at most \(1-2\eta\|A\|_{2}\).

    \item \textbf{Stochastic gradients.}  
    If only unbiased stochastic estimators \(\tilde F(z_t)\) with bounded variance are available, the same potential analysis holds in expectation, leading to an \(O(1/\sqrt{T})\) rate for the stochastic OGDA algorithm (mirroring the SGD result in the reference example).  
\end{enumerate}

%%%%%%%%%%%%%%%%%%%%%%%%%%%%%%%%%%%%%%%%%%%%%%%%%%%%%%%%%%%%
\section*{Discussion}
\addcontentsline{toc}{section}{Discussion}
OGDA combines the low computational cost of plain GDA (one gradient per iteration) with the stability benefits of extragradient methods.  The optimistic correction term effectively predicts the next gradient, cancelling part of the rotational component that causes limit cycles in bilinear games.  The rigorous analysis above shows that under the standard smoothness and convex–concave assumptions the method achieves the optimal \(O(1/T)\) rate for deterministic first‑order methods, with a simple stepsize condition \(\eta\le 1/(4L)\).  In practice, a stepsize close to this bound often yields the fastest empirical convergence, while larger stepsizes may lead to divergence.

%%%%%%%%%%%%%%%%%%%%%%%%%%%%%%%%%%%%%%%%%%%%%%%%%%%%%%%%%%%%
\begin{thebibliography}{9}
\bibitem{daskalakis2018}
C.~Daskalakis, I.~Panageas, and A.~Zampetakis. \emph{Training GANs with Optimistic Mirror Descent}. In \textit{Advances in Neural Information Processing Systems}, 2018.

\bibitem{gidel2020}
G.~Gidel, G.~Bertrand, and P.~Vincent. \emph{A Variational Inequality Perspective on Generative Adversarial Networks}. In \textit{International Conference on Machine Learning}, 2020.

\end{thebibliography}

\end{document}